\documentclass[11pt,a4paper]{article}
\usepackage[utf8]{inputenc}
\usepackage[T1]{fontenc}
\usepackage{ctex}
\usepackage{amsmath,amssymb,amsfonts}
\usepackage{algorithm}
\usepackage{algpseudocode}
\usepackage{booktabs}
\usepackage{graphicx}
\usepackage{hyperref}
\usepackage{xcolor}
\usepackage{listings}
\usepackage{geometry}
\usepackage{tikz}
\usetikzlibrary{positioning,arrows.meta,shapes.geometric,fit,calc}
\geometry{margin=2.5cm}

\lstset{
  language=Python,
  basicstyle=\ttfamily\small,
  keywordstyle=\color{blue}\bfseries,
  commentstyle=\color{gray},
  stringstyle=\color{red},
  showstringspaces=false,
  breaklines=true,
  frame=single,
  numbers=left,
  numberstyle=\tiny\color{gray},
}

\title{\textbf{MoDA 技术报告：原理、Triton Kernel 实现细节\\与 LeWM 整合方案}}
\author{MoDA--LeWM Integration Project}
\date{\today}

\begin{document}
\maketitle
\tableofcontents
\newpage

% =========================================================================
\section{引言与动机}
% =========================================================================

\subsection{问题：深层 Transformer 的信息稀释}

在标准 Transformer 中，第 $i$ 层的隐状态 $\mathbf{h}^{(i)}$ 只能通过残差流（residual stream）间接获取浅层信息：
\begin{equation}
  \mathbf{h}^{(i)} = \mathbf{h}^{(i-1)} + \text{Attn}^{(i)}(\mathbf{h}^{(i-1)}) + \text{FFN}^{(i)}(\cdots)
\end{equation}
经过多层叠加后，早期层的细粒度信息（例如第 2 层对某个特定 pattern 的响应）会被后续层的非线性变换逐渐覆盖和稀释。这就是所谓的\textbf{信息稀释}（information dilution）问题。

对于 LeWM 这样的 world model predictor，predictor 深度可达 64 层。第 60 层想要回忆第 2 层"看到了什么"，只能从叠加了 58 层变换后的残差流中间接提取——信息损失不可避免。

\subsection{MoDA 的解法：Mixture-of-Depth Attention}

MoDA（Mixture-of-Depth Attention）的核心思想是：\textbf{让每一层直接"回看"前面所有层的原始 Key/Value}，而不只是依赖残差流。

具体来说，对于第 $i$ 层、位置 $t$ 的 query $\mathbf{q}_t^{(i)}$，它可以 attend 到：
\begin{enumerate}
  \item \textbf{时间维（temporal）}：同一层的 $\mathbf{k}_0^{(i)}, \mathbf{k}_1^{(i)}, \ldots, \mathbf{k}_t^{(i)}$（因果约束）
  \item \textbf{深度维（depth）}：前面所有层在\emph{同一位置 $t$} 的 $\mathbf{k}_t^{(0)}, \mathbf{k}_t^{(1)}, \ldots, \mathbf{k}_t^{(i-1)}$
\end{enumerate}

而且这两组候选\textbf{在同一个 softmax 里竞争}——网络自己学习什么时候该看时间历史、什么时候该看浅层记忆。

% =========================================================================
\section{MoDA 核心数学}
% =========================================================================

\subsection{符号定义}

\begin{table}[h]
\centering
\begin{tabular}{cl}
\toprule
符号 & 含义 \\
\midrule
$B$ & batch size \\
$T$ & 序列长度（时间步数） \\
$L$ & 当前层之前已缓存的层数（depth cache 的层数） \\
$H$ & attention head 数 \\
$D$ & 每个 head 的维度（dim\_head） \\
$\mathbf{q}, \mathbf{k}, \mathbf{v} \in \mathbb{R}^{B \times T \times H \times D}$ & 当前层的 Query, Key, Value \\
$\mathbf{k}_d, \mathbf{v}_d \in \mathbb{R}^{B \times (T \cdot L) \times H \times D}$ & depth cache（前层 Key/Value） \\
$s$ & scale factor $= 1/\sqrt{D}$ \\
\bottomrule
\end{tabular}
\caption{MoDA 符号表}
\end{table}

\subsection{Unified Softmax Attention}

对于 batch $b$、head $h$、query 位置 $t$，MoDA 的计算过程如下：

\paragraph{Step 1: 时间维 logits（causal）}
\begin{equation}
  \ell_{\text{seq}}[j] = \frac{\mathbf{q}_t \cdot \mathbf{k}_j}{\sqrt{D}}, \quad j = 0, 1, \ldots, t
  \label{eq:logits_seq}
\end{equation}
共 $(t+1)$ 个 logits。因果约束自然满足——只 attend 到 $j \leq t$。

\paragraph{Step 2: 深度维 logits}
\begin{equation}
  \ell_{\text{depth}}[l] = \frac{\mathbf{q}_t \cdot \mathbf{k}_{d}[t, l]}{\sqrt{D}}, \quad l = 0, 1, \ldots, L-1
  \label{eq:logits_depth}
\end{equation}
共 $L$ 个 logits。注意这里 $\mathbf{k}_{d}[t, l]$ 是第 $l$ 层在\emph{同一位置 $t$} 产生的 Key。

\paragraph{Step 3: Unified Softmax}
将两组 logits 拼接后做\emph{一次} softmax：
\begin{equation}
  \boldsymbol{\ell} = \text{concat}(\ell_{\text{seq}},\; \ell_{\text{depth}}) \in \mathbb{R}^{(t+1+L)}
\end{equation}
\begin{equation}
  \mathbf{w} = \text{softmax}(\boldsymbol{\ell}) \in \mathbb{R}^{(t+1+L)}
  \label{eq:unified_softmax}
\end{equation}

\paragraph{Step 4: 加权求和}
\begin{equation}
  \mathbf{o}_t = \underbrace{\sum_{j=0}^{t} w_j \cdot \mathbf{v}_j}_{\text{temporal context}} + \underbrace{\sum_{l=0}^{L-1} w_{t+1+l} \cdot \mathbf{v}_{d}[t, l]}_{\text{depth context}}
  \label{eq:output}
\end{equation}

\subsection{为什么要 Unified Softmax？}

\begin{table}[h]
\centering
\begin{tabular}{lll}
\toprule
方案 & 做法 & 特点 \\
\midrule
独立 attention & 分别对时间和深度做两个 softmax，再加权 & 需要额外混合参数 \\
加法融合 & 分别算 output 再相加 & 无竞争，信息冗余 \\
\textbf{Unified softmax (MoDA)} & 拼接后做一个 softmax & \textbf{自动竞争，零额外参数} \\
\bottomrule
\end{tabular}
\caption{三种深度融合方案对比}
\end{table}

Unified softmax 的核心优势：由于 softmax 的归一化性质 $\sum w_i = 1$，时间和深度的权重\textbf{自动此消彼长}。如果某个 depth key 特别相关，它的 score 高，softmax 就会给它更多权重，同时自动压低时间维的权重。\textbf{不需要任何额外参数来控制混合比例。}

\subsection{LogSumExp (LSE) 的作用}

在 Triton kernel 的实现中，除了输出 $\mathbf{o}$ 之外，还会计算 LSE：
\begin{equation}
  \text{LSE}_t = \log \sum_{i} \exp(\ell_i)
\end{equation}
LSE 在 backward pass 中用于高效重建 softmax 权重 $w_i = \exp(\ell_i - \text{LSE})$，避免重新计算完整的 softmax。这是 Flash Attention 系列 kernel 的标准技巧。

% =========================================================================
\section{Depth Cache：构建与布局}
% =========================================================================

\subsection{逐层累积}

Depth cache 在 Transformer 的 forward pass 中逐层构建：

\begin{algorithm}[H]
\caption{Depth Cache 累积过程}
\begin{algorithmic}[1]
\State $\texttt{k\_cache\_list} \gets []$, $\texttt{v\_cache\_list} \gets []$
\For{$i = 0$ \textbf{to} $\text{depth}-1$}
  \If{$i \geq \texttt{depth\_start\_layer}$ \textbf{and} $\texttt{k\_cache\_list} \neq \emptyset$}
    \State $\mathbf{k}_d, \mathbf{v}_d \gets \texttt{build\_depth\_cache}(\texttt{k\_cache\_list}, \texttt{v\_cache\_list}, T)$
  \Else
    \State $\mathbf{k}_d, \mathbf{v}_d \gets \text{None}$
  \EndIf
  \State $\mathbf{h}^{(i)}, \mathbf{k}^{(i)}, \mathbf{v}^{(i)} \gets \texttt{MoDABlock}_i(\mathbf{h}^{(i-1)}, \mathbf{k}_d, \mathbf{v}_d)$
  \State $\texttt{k\_cache\_list}.\text{append}(\mathbf{k}^{(i)})$
  \State $\texttt{v\_cache\_list}.\text{append}(\mathbf{v}^{(i)})$
\EndFor
\end{algorithmic}
\end{algorithm}

\begin{itemize}
  \item \textbf{Layer 0}：没有前层，$\mathbf{k}_d = \text{None}$ → 退化为纯 causal attention
  \item \textbf{Layer $i$ ($i \geq 1$)}：拿到 layers $0, 1, \ldots, i-1$ 的所有 K/V → $L = i$
  \item \textbf{Layer 63}：depth cache 包含 63 层的 K/V，$L = 63$
\end{itemize}

\subsection{Position-Major 布局}

MoDA kernel 要求 depth cache 按\textbf{position-major} 排列：

\begin{equation}
  \mathbf{k}_d = [\underbrace{k_0^{(0)}, k_0^{(1)}, \ldots, k_0^{(L-1)}}_{\text{position 0 的所有层}},\; \underbrace{k_1^{(0)}, k_1^{(1)}, \ldots, k_1^{(L-1)}}_{\text{position 1 的所有层}},\; \ldots]
\end{equation}

最终 shape: $[B, T \times L, H, D]$。

\paragraph{为什么用 position-major？} Triton kernel 在处理 query 位置 $t$ 时，需要读取该位置所有层的 depth K/V。Position-major 布局保证这些数据在内存中连续，最大化 GPU 的内存带宽利用率。

\paragraph{构建代码对应：}
\begin{lstlisting}
# k_list: list of (B, T, H, D), length L
stacked_k = torch.stack(k_list, dim=2)  # (B, T, L, H, D)
cached_k = stacked_k.reshape(B, T*L, H, D)  # position-major
\end{lstlisting}

\texttt{torch.stack(dim=2)} 把 layer 维度插在 $T$ 和 $H$ 之间，然后 \texttt{reshape} 把 $T$ 和 $L$ 两个维度展平。由于 PyTorch 的 row-major 内存布局，展平后自然是 position-major：先遍历 $L$（内层循环），再遍历 $T$（外层循环）。

% =========================================================================
\section{Triton Kernel 详解}
% =========================================================================

Triton 是 OpenAI 开发的 GPU 编程语言，允许用 Python 风格的代码编写高性能 GPU kernel。MoDA 的 Triton kernel 实现了 Flash Attention 风格的 tiled attention，同时融合了 depth cache 的处理。

\subsection{什么是 Triton？}

\begin{itemize}
  \item \textbf{定位}：介于 CUDA 和 PyTorch 之间的抽象层。比 CUDA 易写，比 PyTorch 的 \texttt{torch.nn.functional} 更底层、更灵活。
  \item \textbf{核心概念}：程序员定义一个 kernel 函数，该函数在 GPU 的多个\textbf{程序实例}（program instances）上并行执行。每个实例处理输出的一个 tile（小块）。
  \item \textbf{优势}：自动处理 GPU 的 shared memory、warp 调度等底层细节，同时提供接近手写 CUDA 的性能。
\end{itemize}

\subsection{Forward Kernel 架构}

MoDA forward kernel (\texttt{parallel\_moda\_fwd\_kernel}) 的并行策略：

\begin{equation}
  \text{grid} = (\underbrace{N_V}_{\text{value 分块数}},\; \underbrace{N_T}_{\text{query 时间分块数}},\; \underbrace{B \times H_Q}_{\text{batch} \times \text{head}})
\end{equation}

每个 kernel 实例负责计算输出的一个 $[\text{BT}, \text{BV}]$ 大小的 tile。

\paragraph{关键常量（tile sizes）：}
\begin{table}[h]
\centering
\begin{tabular}{cll}
\toprule
符号 & 含义 & 典型值 \\
\midrule
\texttt{BT} & query block size（时间维分块） & 64--128 \\
\texttt{BS} & key block size（扫描步长） & 32--64 \\
\texttt{BK} & key dimension block & 64--256 \\
\texttt{BV} & value dimension block & 64--256 \\
\bottomrule
\end{tabular}
\caption{Triton kernel tile 参数}
\end{table}

\subsection{Forward Kernel 的三阶段计算}

\begin{figure}[h]
\centering
\begin{tikzpicture}[
  block/.style={draw, rounded corners, fill=blue!10, minimum width=5.5cm, minimum height=1.2cm, align=center},
  arrow/.style={-{Stealth[length=3mm]}, thick},
  node distance=1.5cm
]
\node[block] (s1) {\textbf{Stage 1: Full Causal Blocks}\\扫描 $k[0 \ldots \text{base\_t\_block\_start})$\\所有 key 都满足因果性};
\node[block, below=of s1] (s2) {\textbf{Stage 2: Partial Causal Block}\\扫描 $k[\text{base\_t\_block\_start} \ldots \text{base\_t\_block\_end})$\\需要逐元素因果 mask};
\node[block, below=of s2] (s3) {\textbf{Stage 3: Depth Blocks}\\扫描 $k_d[\text{depth\_block\_start} \ldots \text{depth\_block\_end})$\\需要 position-match mask};
\draw[arrow] (s1) -- (s2);
\draw[arrow] (s2) -- (s3);
\end{tikzpicture}
\caption{Forward kernel 的三阶段处理流程}
\end{figure}

\subsubsection{Stage 1: Full Causal Blocks（完全因果区域）}

对应源码第 174--220 行。

对于 query block 起始位置 \texttt{q\_block\_start}，所有 key 位置 $j < \text{base\_t\_block\_start}$ 必然满足因果条件 $j \leq t$，因此\textbf{不需要逐元素 mask}。

\begin{lstlisting}[title=Stage 1 核心逻辑（伪代码）]
for i_s in range(0, base_t_block_start, BS):
    b_k = load_block(k, i_s)        # (BK, BS)
    b_v = load_block(v, i_s)        # (BS, BV)
    b_s = dot(b_q, b_k) * scale     # (BT, BS) attention scores
    # Online softmax update (Flash Attention style):
    b_m_new = max(b_m, max(b_s))    # running max
    b_r = exp2(b_m_old - b_m_new)   # rescale factor
    b_p = exp2(b_s - b_m_new)       # new weights
    b_acc = b_acc * b_r + sum(b_p)  # running sum
    b_o = b_o * b_r + dot(b_p, b_v) # running output
\end{lstlisting}

\subsubsection{Stage 2: Partial Causal Block（边界因果区域）}

对应源码第 225--267 行。

当 key 位置和 query 位置在同一个 block 范围内时，需要逐元素检查因果条件：
\begin{equation}
  \text{mask}[i, j] = \left(\lfloor \texttt{o\_q}[i] / G \rfloor \geq \texttt{o\_k}[j]\right)
\end{equation}
其中 $G$ = \texttt{moda\_group\_num}。不满足条件的位置设为 $-\infty$。

\subsubsection{Stage 3: Depth Blocks（深度缓存区域）}

对应源码第 269--337 行。这是 MoDA 的核心创新。

\begin{lstlisting}[title=Stage 3 核心逻辑（伪代码）]
for depth_col_start in range(depth_block_start, depth_block_end, BS):
    b_k_depth = load_block(cached_k, depth_col_start)
    b_v_depth = load_block(cached_v, depth_col_start)
    b_s_depth = dot(b_q, b_k_depth) * scale

    # Position-match mask:
    depth_col_ids = depth_col_start + arange(BS)
    depth_row_ids = depth_col_ids // L    # which position
    mask = (query_position == depth_row_ids)  # same position only!
    b_s_depth = where(mask, b_s_depth, -inf)

    # Continue online softmax (shared with Stage 1 & 2):
    b_m_new = max(b_m, max(b_s_depth))
    b_r = exp2(b_m - b_m_new)
    b_p = exp2(b_s_depth - b_m_new)
    b_acc = b_acc * b_r + sum(b_p)
    b_o = b_o * b_r + dot(b_p, b_v_depth)
\end{lstlisting}

\paragraph{关键 mask 逻辑：}
\begin{equation}
  \texttt{depth\_row\_ids} = \lfloor \texttt{depth\_col\_ids} / L \rfloor
\end{equation}
这利用了 position-major 布局：在 $\mathbf{k}_d$ 的线性索引中，连续的 $L$ 个元素属于同一个位置。因此 $\lfloor \text{index} / L \rfloor$ 就是对应的时间位置。只有 query 位置和 depth key 位置匹配时，attention 才有效。

\paragraph{Online Softmax 的统一性：}
三个 stage 共享同一组 running statistics ($\texttt{b\_m}$, $\texttt{b\_acc}$, $\texttt{b\_o}$)。这意味着 temporal scores 和 depth scores \textbf{在同一个 softmax 归一化下}竞争——这正是 MoDA "unified softmax" 的实现方式。

最终归一化：
\begin{equation}
  \texttt{b\_o} = \texttt{b\_o} / \texttt{b\_acc}
\end{equation}

\subsection{Online Softmax（Flash Attention 技巧）}

标准 softmax 需要两遍扫描（第一遍求 max，第二遍求 exp/sum），这对 GPU 不友好。Online softmax 用一遍扫描完成，核心是维护三个 running statistics：

\begin{align}
  m_{\text{new}} &= \max(m_{\text{old}}, \max(\mathbf{s}_{\text{new}})) \label{eq:online_m} \\
  r &= 2^{m_{\text{old}} - m_{\text{new}}} \quad \text{（rescale factor）} \label{eq:online_r} \\
  \text{acc}_{\text{new}} &= \text{acc}_{\text{old}} \cdot r + \sum_j 2^{s_j - m_{\text{new}}} \label{eq:online_acc} \\
  \mathbf{o}_{\text{new}} &= \mathbf{o}_{\text{old}} \cdot r + \sum_j 2^{s_j - m_{\text{new}}} \cdot \mathbf{v}_j \label{eq:online_o}
\end{align}

注意这里用的是 $2^x$（\texttt{exp2}）而不是 $e^x$，因此 scores 预乘了 $1/\ln 2$（即 \texttt{RCP\_LN2 = 1.4427}）。这是因为 GPU 硬件上 $2^x$ 比 $e^x$ 快。

\subsection{Backward Kernel 架构}

Backward pass 分为\textbf{四个独立 kernel}：

\begin{enumerate}
  \item \texttt{parallel\_moda\_bwd\_kernel\_preprocess}：计算 $\delta = \sum_d o_d \cdot do_d$（逐元素 dot product）
  \item \texttt{parallel\_moda\_bwd\_kernel\_dq}：计算 $\frac{\partial \mathcal{L}}{\partial \mathbf{q}}$，包含时间维和深度维
  \item \texttt{parallel\_attn\_bwd\_kernel\_dkv}：计算时间维的 $\frac{\partial \mathcal{L}}{\partial \mathbf{k}}$, $\frac{\partial \mathcal{L}}{\partial \mathbf{v}}$
  \item \texttt{parallel\_attn\_bwd\_kernel\_dkv\_depth}：计算深度维的 $\frac{\partial \mathcal{L}}{\partial \mathbf{k}_d}$, $\frac{\partial \mathcal{L}}{\partial \mathbf{v}_d}$
\end{enumerate}

\subsubsection{Backward 的核心公式}

利用前向保存的 LSE，后向可以高效重建 softmax 权重：
\begin{equation}
  p_{ij} = 2^{s_{ij} - \text{LSE}_i}
\end{equation}

然后梯度计算：
\begin{align}
  dp_{ij} &= \mathbf{do}_i \cdot \mathbf{v}_j \\
  ds_{ij} &= p_{ij} \cdot (dp_{ij} - \delta_i) \quad \text{where } \delta_i = \sum_j p_{ij} \cdot dp_{ij} = \mathbf{o}_i \cdot \mathbf{do}_i \\
  d\mathbf{q}_i &= \sum_j ds_{ij} \cdot \mathbf{k}_j \cdot s \\
  d\mathbf{k}_j &= \sum_i ds_{ij} \cdot \mathbf{q}_i \cdot s \\
  d\mathbf{v}_j &= \sum_i p_{ij} \cdot \mathbf{do}_i
\end{align}

对于 depth cache 的梯度 $d\mathbf{k}_d$, $d\mathbf{v}_d$，计算方式完全类似，只是 mask 从因果 mask 变为 position-match mask。

\subsection{GPU 适配与 Tile 参数}

Kernel 会根据 GPU 型号选择不同的 tile 参数：

\begin{table}[h]
\centering
\begin{tabular}{lccccl}
\toprule
GPU & BT & BS & BK & BV & warps \\
\midrule
H800 (Hopper) & 64 & 64 & 64 & 64 & 4 \\
H20 & 64 & 64 & 64 & 64 & 4 \\
A100 (Ampere) & 64 & 64 & 64 & 64 & 4 \\
\textbf{其他（如 RTX 6000 Ada）} & 128 & 32 & 256 & 64 & 2 \\
\bottomrule
\end{tabular}
\caption{不同 GPU 的 forward kernel tile 参数}
\end{table}

RTX 6000 Ada 不在优化列表中，使用 fallback 参数。\texttt{Target GPU is not optimized} 的 warning 来自这里。\textbf{数学上完全正确}，只是 tile 大小不是针对该 GPU 手动调优的。

\subsection{autograd 整合}

MoDA 通过 PyTorch 的 \texttt{torch.autograd.Function} 将 Triton kernel 接入自动微分系统：

\begin{lstlisting}[title=autograd 接口（简化）]
class ParallelMoDAFunction(torch.autograd.Function):
    @staticmethod
    def forward(ctx, q, k, v, cached_k, cached_v, scale, ...):
        o, lse = parallel_moda_fwd(q, k, v, ...)
        ctx.save_for_backward(q, k, v, o, lse, cached_k, cached_v)
        return o

    @staticmethod
    def backward(ctx, do):
        q, k, v, o, lse, cached_k, cached_v = ctx.saved_tensors
        dq, dk, dv, _, d_cached_k, d_cached_v = parallel_moda_bwd(
            q, k, v, o, ..., do, cached_k=cached_k, cached_v=cached_v
        )
        return dq, dk, dv, d_cached_k, d_cached_v, None, ...
\end{lstlisting}

forward 保存 $\mathbf{o}$ 和 LSE（而不是完整的 attention weights），backward 利用 LSE 重建权重。这是 Flash Attention 的核心 memory saving 技巧。

% =========================================================================
\section{MoDA 额外特性}
% =========================================================================

\subsection{moda\_group\_num}

当 \texttt{moda\_group\_num} $= G > 1$ 时，连续 $G$ 个 query token 共享同一个 base 位置：
\begin{equation}
  \text{base\_t} = \lfloor t_q / G \rfloor
\end{equation}

这允许多个 query 映射到同一个"物理时间步"的 depth cache。LeWM 中 $T=3$，不需要此特性，固定 $G=1$。

\subsection{Gate $g$（cumsum gate）}

MoDA 支持可选的位置衰减 gate $g$。启用时，attention score 会加上一个 cumulative sum bias：
\begin{equation}
  s'_{ij} = s_{ij} + g_{\text{cumsum}}[i] - g_{\text{cumsum}}[j]
\end{equation}
这给远距离的 key 施加衰减。LeWM 不使用此特性（$g = \text{None}$）。

\subsection{GQA（Grouped-Query Attention）}

MoDA 支持 $H_Q \neq H$（多个 query head 共享同一组 K/V head）。LeWM 使用 $H_Q = H = 16$，不涉及 GQA。

% =========================================================================
\section{LeWM 整合方案}
% =========================================================================

\subsection{架构对比}

\begin{figure}[h]
\centering
\begin{tikzpicture}[
  block/.style={draw, rounded corners, fill=green!10, minimum width=4cm, minimum height=0.8cm, align=center, font=\small},
  moda/.style={draw, rounded corners, fill=orange!20, minimum width=4cm, minimum height=0.8cm, align=center, font=\small},
  arrow/.style={-{Stealth[length=2mm]}, thick},
  node distance=0.8cm
]
% Left: Original LeWM
\node[font=\bfseries] at (-3, 5) {原版 LeWM};
\node[block] (l1) at (-3, 4) {AdaLN-zero 调制};
\node[block] (l2) at (-3, 3) {标准 Causal Attn};
\node[block] (l3) at (-3, 2) {FFN};
\node[block] (l4) at (-3, 1) {残差连接};
\draw[arrow] (l1) -- (l2);
\draw[arrow] (l2) -- (l3);
\draw[arrow] (l3) -- (l4);

% Right: MoDA LeWM
\node[font=\bfseries] at (3, 5) {MoDA 版 LeWM};
\node[block] (r1) at (3, 4) {AdaLN-zero 调制};
\node[moda] (r2) at (3, 3) {\textbf{MoDA Unified Attn}};
\node[block] (r3) at (3, 2) {FFN};
\node[block] (r4) at (3, 1) {残差连接};
\draw[arrow] (r1) -- (r2);
\draw[arrow] (r2) -- (r3);
\draw[arrow] (r3) -- (r4);

% Arrow between
\draw[-{Stealth[length=3mm]}, ultra thick, red] (l2) -- node[above, font=\small\bfseries, red] {唯一改动} (r2);
\end{tikzpicture}
\caption{原版 LeWM 与 MoDA 版的唯一区别：attention 模块}
\end{figure}

\subsection{保留的 LeWM 组件}

\begin{itemize}
  \item \textbf{AdaLN-zero 调制}：$\text{modulate}(\mathbf{x}, \text{shift}, \text{scale}) = \mathbf{x} \cdot (1 + \text{scale}) + \text{shift}$，action embedding 通过 6 路 AdaLN 调制 attention 和 FFN
  \item \textbf{ViT-tiny encoder}：hidden\_size=192, 12 层, 3 heads
  \item \textbf{SIGReg loss}：$\mathcal{L} = \mathcal{L}_{\text{pred}} + \lambda \cdot \mathcal{L}_{\text{sigreg}}$
  \item \textbf{Learnable positional embedding}（而非 MoDA 原版的 RoPE）
  \item \textbf{无 QK Norm}（MoDA 原版有）
\end{itemize}

\subsection{MoDA 版 LeWM 的 Attention 流程}

\begin{algorithm}[H]
\caption{MoDAAttention.forward()}
\begin{algorithmic}[1]
\Require $\mathbf{x} \in \mathbb{R}^{B \times T \times D}$（经 AdaLN-zero 调制后）, $\mathbf{k}_d$, $\mathbf{v}_d$（depth cache 或 None）
\State $\mathbf{x}_{\text{normed}} \gets \text{LayerNorm}(\mathbf{x})$
\State $\mathbf{q}, \mathbf{k}, \mathbf{v} \gets \text{split}(\text{Linear}(\mathbf{x}_{\text{normed}}))$ \Comment{各 $\in \mathbb{R}^{B \times T \times H \times D}$}
\State $\mathbf{o} \gets \texttt{parallel\_moda}(\mathbf{q}, \mathbf{k}, \mathbf{v}, \text{cached\_k}=\mathbf{k}_d, \text{cached\_v}=\mathbf{v}_d, \text{scale}=1/\sqrt{D})$
\State $\mathbf{o} \gets \text{Linear}(\text{reshape}(\mathbf{o}))$
\State \Return $\mathbf{o}$, $\mathbf{k}$, $\mathbf{v}$ \Comment{$\mathbf{k}, \mathbf{v}$ 存入 depth cache}
\end{algorithmic}
\end{algorithm}

\subsection{配置参数}

\begin{table}[h]
\centering
\begin{tabular}{ll}
\toprule
参数 & 值 \\
\midrule
\texttt{use\_moda} & \texttt{true} \\
\texttt{depth\_start\_layer} & 1（第 0 层无 depth cache） \\
\texttt{predictor.depth} & 64 \\
\texttt{predictor.heads} & 16 \\
\texttt{predictor.dim\_head} & 64 \\
\texttt{predictor.mlp\_dim} & 2048 \\
\texttt{wm.history\_size} & 3（序列长度 $T=3$） \\
\texttt{trainer.precision} & bf16 \\
\bottomrule
\end{tabular}
\caption{LeWM MoDA 实验默认配置}
\end{table}

% =========================================================================
\section{实现一致性验证}
% =========================================================================

\subsection{三层验证体系}

\begin{table}[h]
\centering
\begin{tabular}{llc}
\toprule
验证项 & 方法 & 结果 \\
\midrule
LeWM wrapper $=$ 官方 kernel（有 depth） & \texttt{wrapped - direct = 0} & max\_abs $= 0.0$ \\
LeWM wrapper $=$ 官方 kernel（无 depth） & \texttt{wrapped - direct = 0} & max\_abs $= 0.0$ \\
官方 kernel $\approx$ 官方 naive reference & bf16 kernel 容差内 & max\_abs $< 0.01$ \\
\bottomrule
\end{tabular}
\caption{数值正确性验证结果}
\end{table}

第三项的误差来自 bf16 精度下 Triton kernel 的 tiled 计算与 Python for-loop 的 fp32 参考实现之间的数值差异，属于 kernel 正常容差范围。

\subsection{与原版 MoDA 的逐项对比}

\begin{table}[h]
\centering
\small
\begin{tabular}{p{2.8cm}p{4cm}p{4cm}c}
\toprule
维度 & MoDA 原版 & 当前 LeWM 实现 & 一致性 \\
\midrule
核心数学 & unified softmax over [seq, depth] & 直接调用官方 \texttt{parallel\_moda} & \checkmark \\
depth cache 格式 & $[B, T{\cdot}L, H, D]$ position-major & \texttt{stack(dim=2).reshape} & \checkmark \\
depth cache 来源 & 前面所有层的 K/V & \texttt{k\_cache\_list} 累积 & \checkmark \\
depth 开始层 & 第 1 层开始 & \texttt{depth\_start\_layer=1} & \checkmark \\
causal masking & kernel 内部实现 & 直接走官方 kernel & \checkmark \\
scale & $1/\sqrt{D}$ & \texttt{dim\_head**-0.5} & \checkmark \\
Triton 加速 & 官方 \texttt{parallel\_moda} & 同一份代码 & \checkmark \\
\midrule
\texttt{moda\_group\_num} & 支持 $>1$ & 固定 $=1$ & 简化 \\
gate $g$ & 支持 cumsum gate & 固定 None & 简化 \\
GQA & 支持 $H_Q \neq H$ & $H_Q = H$ & 简化 \\
Transformer block & LLaMA-like & AdaLN-zero (LeWM) & 刻意保留 \\
位置编码 & RoPE & learnable & 刻意保留 \\
QK Norm & 有 & 无 & 刻意保留 \\
\bottomrule
\end{tabular}
\caption{完整对比表}
\end{table}

% =========================================================================
\section{总结}
% =========================================================================

\begin{enumerate}
  \item \textbf{MoDA 的核心}是 unified softmax——让时间维和深度维的 K/V 在同一个 softmax 里竞争，网络自动学习注意力分配。
  \item \textbf{Triton kernel} 实现了 Flash Attention 风格的 tiled online softmax，分三阶段处理：完全因果块 → 边界因果块 → 深度缓存块。三个阶段共享 running statistics，实现 unified softmax。
  \item \textbf{Depth cache} 按 position-major 布局，保证 GPU 内存访问连续。
  \item \textbf{Backward} 分四个 kernel：预处理 $\delta$、计算 $d\mathbf{q}$、计算时间维 $d\mathbf{k}/d\mathbf{v}$、计算深度维 $d\mathbf{k}_d/d\mathbf{v}_d$。利用 LSE 重建 softmax 权重，避免保存完整 attention matrix。
  \item \textbf{LeWM 整合}只替换了 attention 模块（标准 causal → MoDA unified），其他全部保留，确保实验结果可归因于 depth attention 机制。
\end{enumerate}

\end{document}
